\chapter{2010 Nobel Prize in Economics}
\textbf{Laureates: }Peter Mortensen (USA/Denmark), Christopher Pissarides (UK/Cyprus)

\section*{Reason for Award}
For their analysis of markets with search frictions

\section{What They Did}
Dale Mortensen and Christopher Pissarides developed search and matching theory,
explaining how labor markets function despite frictions that prevent instant
worker-firm matches. Their model showed that unemployment persists because
finding suitable matches takes time and resources.

Job seekers search for vacancies, firms search for workers, and the matching
process determines the unemployment rate. Mortensen contributed the equilibrium
wage-setting model, showing how bargaining between workers and firms determines
wages. He analyzed how unemployment benefits affect search intensity and job
acceptance rates.

Pissarides developed the matching function approach, quantifying how vacancies
and unemployment combine to produce hires. Together, they created a framework
for analyzing unemployment dynamics, job creation and destruction, and the
effectiveness of labor market policies.

Their model explained why unemployment can persist even when wages are flexible,
due to coordination failures in the matching process.

\section{Impact on Economic Thinking}
Search and matching theory transformed labor economics by providing microfoundations
for unemployment. Before their work, unemployment was often explained by rigid
wages or exogenous shocks. The search framework showed that even with flexible
wages, unemployment can persist due to coordination failures in the matching
process.

This gave policymakers new tools for analyzing unemployment benefits, job search
requirements, and active labor market programs. The theory explained why
unemployment insurance affects search intensity and job acceptance rates.

It also illuminated the determinants of vacancy creation and the efficiency of
labor market institutions. Business cycle fluctuations in unemployment could now
be understood through changes in matching efficiency and job destruction rates.

The framework influenced macroeconomic modeling, with search frictions incorporated
into dynamic stochastic general equilibrium models, creating the DSGE approach
to business cycle analysis.

\section{Modern Perspectives Today}
Search and matching theory remains central to labor economics and macroeconomics.
Modern applications analyze the Great Recession's unemployment dynamics, the gig
economy's impact on job search, and the effects of automation on matching
efficiency.

Policymakers use matching models to evaluate unemployment benefit reforms, job
training programs, and employment subsidies. The framework has been extended to
analyze housing markets, credit markets, and platform economies where matching
frictions are prevalent.

Digital labor platforms like Uber and Upwork have created new matching
technologies that economists study using search theory. Research on wage
inequality, job mobility, and career dynamics builds on the Mortensen-Pissarides
framework.

The theory also informs debates about minimum wage policies, unionization, and
worker protections. As labor markets evolve with technology and globalization,
search and matching theory provides essential analytical tools.
